% !TeX root = ../papers/nc-p3.tex
\section{分岐方向の二つの平方根と接線の剰余座標}
\label{sec:rank-fifteen-ramified-pole}
\begin{lemma}\label{lem:rank-fifteen-ramified-pole}
剰余標数3で $\bar x=1$, $\nu(y)>0$, $2\nu(x-1)\ge\nu(y)$ とする。
$J=(x-1)^2/(4y)$ は整である。付値を延長すると、$\nu(H)=\nu(y)/4$ と
単元 $z$ による表示があり、$\bar J=(z^4+1)^2/(4z^4)\bmod\mathfrak m$ となる。
$\bar J\ne1$ ならこの式のどの幾何学的剰余根も指定できる。
\end{lemma}
\begin{proof}
前節と同じ $U=1+r^2$, $K=\sqrt{1+r^2}$, $x=(4+r^2)T^2$,
$y=r^4T^6$ において $K\equiv1$, $2T\equiv1$ を選ぶ。
$h=r/(K+1)$, $H=h/\sqrt2$ とし、$v^2=1-2T$, $u^2=1+2T$ と置く。
恒等式 $1-2T=(1+r^2/4-x)/((1+r^2/4)(1+2T))$ と仮定から
$L=v/h$ は整である。第二の小さい平方根の中身は
\[
  2/(K+1)-2T=v^2-h^2.
\]
$C^2=L^2-1$ と取ると $(L+C)(L-C)=1$ なので $z=L+C$ は単元である。
$L=(z+z^{-1})/2$, $C=(z-z^{-1})/2$ によりこの平方根を $hC$ と書ける。
一方 $A=(u-v)/2$, $B=(u+v)/2$ とすれば $A^2+B^2=1$, $AB=T$ である。

$D=(1-2H^2)^{-1}$, $E=\sqrt{1-H^2 L^2}$, $F=\sqrt{1-H^2(L^2+1)}$ の
平方根の定数項を1とする。$K=(1+2H^2)D$, $r^2=8H^2D^2$ となり、
経路の係数に入る量は
\[
 \begin{gathered}
  rA=2HD(E-HL),\quad rB=2HD(E+HL),\quad T=(1-2H^2L^2)/2,\\
  P_\pm=\sqrt{4D\pm 2Tr^2},\qquad Q_+=4HDF,\quad Q_-=4H^2DC .
 \end{gathered}
\]
$P_\pm$ の定数項は2である。従って係数は $2z$ を逆にした環上で展開できる。
$x=(1+r^2/4)(1-h^2 L^2)^2$ と $h^4/(4y)\equiv1$ より
$(x-1)/h^2\equiv1-2L^2=-(z^4+1)/(2z^2)$ であり、所要の剰余式を得る。

$r^2$ を与える二つの単純な解は第\ref{sec:rank-fifteen-smith-slopes}節の
$(4+de)^3=e^2$ の $e\equiv\pm8$ に対応し、それらの $h^2$ の先頭は互いに逆符号である。
従って $(x-1)/h^2$ の剰余の二つの符号を選べる。
一方、$v,C$ の符号により $z,-z,1/z,-1/z$ を選べる。
残る符号の交換は $z\mapsto iz$, $i^2=-1$ に対応する。
$\bar J=0$ では $z^4=-1$ の四根は既に最初の四表示で覆われる。
$\bar J\ne0,1$ では八表示は相異なる。
\end{proof}

\begin{theorem}\label{thm:rank-fifteen-ramified-pole}
補題の条件の下、通常な階数15窓で $\bar J\ne1$,
$0<3\nu(y)\le2\nu(3)$ とする。この窓は標準Koszul左右AS正則であり、
各始点で四中心四次式上自由階数256である。特に $x=1$ の全ての通常点で
$0<3\nu(y)\le2\nu(3)$ ならこの結論が成り立つ。
\end{theorem}
\begin{proof}
\path{computations/verify_rank_fifteen_ramified_pole.py} はまず
$\mathbb F_3(z)$ 上で $H$ の十七桁を使い、$m=z^2+z+2$ で単元の主元を選ぶ。
未切断経路の長さ1から5では最大 $(0,2,2,2,2)$ 桁を使い、中心式で一桁、
五次の積表で一桁を使う。長さ6から13での最大損失は $(1,0,1,1,2,0,1,0)$
であり、十三次で一桁を残して商が消える。

逆元に用いた定数係数の分子・分母の平方自由な積は次数897である。
$z=0,\infty$ を外して、$z\mapsto-z,iz,1/z$ の引き戻しで共通の零点を取る。
$-z$ の処理後は偶多項式なので $iz$ の代入にも $\mathbb F_3$ 上の計算で足りる。
得られる次数100の平方自由式は、$z^4-1$ と次の五つの多項式の
$J=(z^4+1)^2/(4z^4)$ によるファイバーの積である（分母と定数倍を除く）。
同じ列に、所定のファイバーで使う $z$ の既約式を記す。
\[
\begin{array}{c|l}
 J+1 & z^4+z^2+2\\
 J^4+J^3+J^2+1 & z^4+z+2\\
 J^2+1 & z^4+z^3+z^2+z+1\\
 J^2+J+2 & z^4+2z^3+z^2+1\\
 J^3+J^2+2 & z^6+2z^5+2z^3+2z+1
\end{array}
\]
$z^4-1$ は $J=1$ のファイバーの被約部分である。
他の五つの代表で同時に単元となる主元を選んで同じ有理関数上の消去を行う。
未切断経路80組、中心4組、五次までの中心商96組、左右の積の照合128組で、
全308組の主元次数の列は基準表示と一致する。四次中心の自由座標を
指定する語基底と列番号も一致する。これで $\bar J\ne1$ の全ての幾何点に剰余表示がある。

標数零では基準表示で $\mathbb Q(z)$ 上の十二桁を使って長さ7まで進み、
各 $H$ 除算の前に定数係数の恒等的な消滅を検査する。
上記308組のうち最初の212組で標数3と同じ全指数を得る。
補題\ref{lem:rank-fifteen-smith-transfer}により、他の良い表示でもこの範囲の
語の格子と積表は混標数で持ち上がる。実際、二つの積の照合行列で前段の
語基底を替える操作は、左右の生成系と中間語の空間にそれぞれ可逆な基底変更を施す。
中心座標を揃えてあるので同一の商でこの帰納を行える。
七次の剰余積表に六桁が残り、仮定は $3\in H^6\mathcal O$ と同値である。
その六桁から再開した八次から十三次の次元上界は
$(1+t+t^2+t^3)^4$ の全係数に一致するので、前節と同じ自由性と左右の解消の判定が適用できる。
\end{proof}

\begin{theorem}\label{thm:rank-fifteen-collision-tube}
同じ通常窓で $\nu(y)>0$, $\bar x=1$ とする。
$\nu(J-1)\ge9\nu(y)$ なら、深さの上限なしに標準Koszul左右AS正則であり、
各始点で四中心四次式上自由階数256である。$J=1$ も含む。
\end{theorem}
\begin{proof}
ここでは $J$ が単元なので $2\nu(x-1)=\nu(y)$ である。
$r^2$ の二表示のうち $L\equiv0$ となる方を使い、$z\equiv i$, $i^2=-1$ と選ぶ。
$a=H^2$, $l=L^2$, $t=1-2al$, $d=2a/(1-2a)^2$ と置くと
$x=(1+d)t^2$, $2\sqrt y=dt^3$ である。実際の二因子は
\[
\begin{aligned}
(1-2a)^2(x-1+2\sqrt y)&=-4a(l-1)(al-1)(4a^2l-1),\\
(1-2a)^2(x-1-2\sqrt y)&=4al(al+a-1)(4a^2l-4a+1).
\end{aligned}
\]
従って $J-1$ は $L^2(L^2-1)$ の単元倍である。
$2zL=(z-i)(z+i)$ より $\nu(z-i)=\nu(J-1)/2\ge18\nu(H)$。
$J=1$ のときには $z=i$ となる。

補題の式で同じ $H$ を固定して $z$ を $i$ にすると、各経路係数の差は
$H^2(z-i)$ の倍数である。実際、$rA,rB,Q_-$ で最初に $z$ に依存する項は
$H^2 L$ または $H^2 C$ に単元を掛けたもので、残りの単元平方根の差は
根の和で割って求められる。従って窓の係数は $H^{20}$ まで一致する。
コードは $\mathbb Q(i)$ 上でこの二十桁を計算する。全定数項の消滅は厳密であり、
逆元のノルムを3で還元すると非零である。未切断の各長さで $(0,2,2,3,2)$、
中心で一桁、五次の中心商で一桁を使う。長さ6から13での最大損失は
$(1,0,1,1,2,0,3,0)$ で、十三次の後に一桁が残る。
四次の中心商も一桁を使うが、五次の中心商は未切断経路の基底から直接作る。
この有限精度の主元・除算を実際の窓へ移し、中心の一般の下界と合わせて
中心式を持ち上げる。商の次元上界が全て一致するので全次数の結論を得る。
\end{proof}
